%% toc: P15 165 165.5
\chapter{Ranking Theory: Attributed Claims, Contraction and Consonant Belief Functions}
\label{sec:p15-w02b-p15-039}
\index{ranking theory}\index{belief contraction}

Chapter~\ref{sec:p15-rationality-5-3} sets out the ranking theory of the 2021
chapter by Gabriele Kern-Isberner, Niels Skovgaard-Olsen and Wolfgang Spohn:
negative, positive and two-sided ranks, conditional ranks, the reason
relation, and the two-cell update by which a ranking state learns.  The
present chapter continues that treatment in four directions.  It fixes the
finite set of works on which the chapter's claims depend and accounts, claim
by claim, for what its pages 341--343 attribute to those works.  It proves the
Levi and Harper identities for the chapter's two-cell update and shows that
contraction is a case of that update, so that a contracted state can be
changed again.  It states the chapter's notion of independence of one
proposition from another in product form and characterizes it by invariance
under revision.  And it makes exact the chapter's comparison of ranks with
consonant belief functions, separating agreement on conditioning from
divergence under Dempster's rule.  Here ``the chapter'' and ``the source''
always mean the 2021 chapter, and every definition and result labelled in
Chapter~\ref{sec:p15-rationality-5-3} is used exactly as stated there
\parencite[pp.~337--345]{KernIsbernerSkovgaardOlsenSpohn2021Ranking}.

\section{Attributed claims and the works they depend on}
\label{sec:p15-w02b-p15-039-dependents}

Pages 341--343 of the chapter compare ranking theory with other formalisms
and survey work in artificial intelligence and psychology.  Most statements
there are attributions: they report what another work proposed, proved or
found.  An attribution is evidence of the cited work's content only as the
chapter describes it.  To keep such reports apart from what is derived here,
the works on which the chapter's claims depend are fixed first, and every
attributed claim of those pages is then given its kind, its dependent works
and the treatment it receives in the present chapter or in
Chapter~\ref{sec:p15-rationality-5-3}.

\begin{hbconvention}{def:p15-w02b-p15-039-dependent-set}{The chapter's finite dependent-source set}
The dependent-source set \(\mathcal D\) of the chapter is its printed
reference list, pp.~344--345, which has 48 entries.  Grouped by the page of
the first citing passage they are:
p.~337, McGrath 2012 and Spohn 1988, 2009, 2012, 2015, 2020;
p.~338, Christensen 2004;
p.~339, Carnap 1950/1962;
p.~340, Hild and Spohn 2008, Pearl 1988, Goldszmidt and Pearl 1996,
Spohn 2018;
p.~341, Shackle 1949, Rescher 1964, Cohen 1970, 1977 and 1980, Zadeh 1978,
Shafer 1976, Leitgeb 2017, Raidl and Skovgaard-Olsen 2017, Makinson 1989,
Pearl 1990;
p.~342, Darwiche and Pearl 1997, Jin and Thielscher 2007, Delgrande and
Jin 2012, Kern-Isberner and Huvermann 2017, Kern-Isberner, Bock, Sauerwald
and Beierle 2017, Kern-Isberner 2004, Baratgin and Politzer 2010, Oaksford
and Chater 2013, Wolf, Rieger and Knauff 2012, Colombo, Elkin and Hartmann
2018, Da Silva Neves, Bonnefon and Raufaste 2002, Benferhat, Bonnefon and
Da Silva Neves 2005, Isberner and Kern-Isberner 2017, Eichhorn,
Kern-Isberner and Ragni 2018;
p.~343, Ragni, Eichhorn, Bock, Kern-Isberner and Tse 2017, Byrne 1989,
Spohn 2013, Skovgaard-Olsen 2014, Skovgaard-Olsen, Singmann and Klauer 2016
and 2017, Skovgaard-Olsen, Kellen, Hahn and Klauer 2019, Henrion, Provan,
Del Favero and Sanders 1994, Juslin, Nilsson, Winman and Lindskog 2011;
p.~344, note~6, Jeffrey 1983 and Shenoy 1991.
Spohn 2012 is cited at two specific places, its chapter~7 (p.~342) and its
section~11.9 (p.~341).  Other chapters of the same handbook to which the
chapter points are cross-references, not members of \(\mathcal D\).

A claim resting on a member of \(\mathcal D\) is reported below as the
chapter's attribution, with the chapter's page, and is not extended beyond
what the chapter prints.  Two members, the 2016 relevance experiment and
the 2011 logarithm-format study, are reported in the psychology section of
Chapter~\ref{sec:p15-rationality-5-3} from
their own pages.  Where an attributed claim has a formal content that can
be stated with the chapter's own definitions, that content is proved here
and the proof uses no member of \(\mathcal D\)
\parencite[pp.~337--345]{KernIsbernerSkovgaardOlsenSpohn2021Ranking}.
\end{hbconvention}

\subsection{The attributions of pages 341--343}
\label{sec:p15-w02b-p15-039-attributions}

\paragraph{Probability and outright belief: the stability proposals (p.~341).}
The chapter says that belief in \(A\) cannot be represented as
\(P(A)=1-\varepsilon\) while belief sets stay consistent and deductively
closed, that the relation between probability and belief is disputed in
philosophy with no clear solution in view, and it names Leitgeb 2017 and
Raidl and Skovgaard-Olsen 2017 as examples of proposals.  This is the
chapter's reason for developing ranking theory as a theory of belief in its
own right.  Kind: the first statement is a formal claim; the second reports
an open philosophical debate; the third is a methodological decision.  The
formal claim is the ten-outcome contrast of Chapter~\ref{sec:p15-rationality-5-3}, where a fixed threshold
below \(0.9\) accepts jointly inconsistent propositions, and
Theorem~\ref{thm:p15-rat53-possibility} shows that ranks are recovered from
probabilities only as limiting exponents of a family, not from one
probability with one acceptance threshold.  The chapter prints neither
stability proposal's definitions nor its theorems, so neither is stated here,
and nothing here decides whether either succeeds.
Chapter~\ref{sec:p15-rationality-4-7} mentions Leitgeb's stability account
as a separately cited presentation
\parencite[pp.~338, 341]{KernIsbernerSkovgaardOlsenSpohn2021Ranking}.

\paragraph{Dempster--Shafer theory (p.~341).}
The chapter makes three claims.  First, the theory of Dempster--Shafer
belief functions of Shafer 1976 is far more general and contains both
probability theory and ranking theory as special cases.  Second, the
consonant belief functions of Shafer's chapter~10 appear to be equivalent to
negative ranking functions.  Third, the dynamic behaviour of the two diverges,
which prevents a reduction of ranking theory to Dempster--Shafer theory; for
this it cites Spohn 2012, section~11.9.  Kind: formal.  The probability case
of the first claim is the singleton-focal-set case of
Theorem~\ref{thm:p15-r47-focal}.  The ranking case and the second claim are
made exact for finite spaces in part~(i) of
Theorem~\ref{thm:p15-w02b-p15-039-consonant} below.  The argument of
Spohn's section is not printed in the chapter and is not reconstructed here.
Parts~(ii) and~(iii) of the same theorem establish one precise sense of the
third claim: conditioning on a proposition gives the same result in both
calculi, while Dempster's combination rule leads out of the class of belief
functions that correspond to ranks.  Whether this is the divergence that
Spohn's section establishes cannot be read off from the chapter
\parencite[p.~341]{KernIsbernerSkovgaardOlsenSpohn2021Ranking}.

\paragraph{The probability parallel and the indirect route (pp.~341, 343).}
On p.~341 the chapter translates sums of probabilities into minima of ranks
and products and quotients into sums and differences; this is treated
in Theorem~\ref{thm:p15-w02b-p15-039-chain} and
Theorem~\ref{thm:p15-rat53-possibility}.  On p.~343 it characterizes the
reason-relation line of psychological research as indirect: that research
tests ideas from ranking theory through their extension to conditionals in
Spohn 2013 and 2015 and through their parallels with probability theory, and
the chapter adds a pointer to Henrion, Provan, Del Favero and Sanders 1994.
Kind: a methodological characterization.  The 1994 paper is a conference
contribution whose title announces an experimental comparison of numerical
and qualitative probabilistic reasoning; the chapter reports no task,
population, design, measure, uncertainty or result from it, and none is
supplied here.  Here it supports no conclusion about human
reasoning or about the relative merit of numerical and qualitative methods
\parencite[pp.~341, 343]{KernIsbernerSkovgaardOlsenSpohn2021Ranking}.

\paragraph{Independence in iterated and multiple revision (p.~342).}
The chapter says that Darwiche and Pearl 1997 proposed general postulates for
iterated revision of epistemic states and illustrated them with ranking
functions, and that Jin and Thielscher 2007 and Delgrande and Jin 2012 did
the same when they proposed new postulates for iterated and for multiple
revision.  It adds that the independence properties proposed in those studies
can be related to independence with respect to ranking functions, for which
it cites Spohn 2012, chapter~7, in analogy with probabilistic independence,
for which it cites Kern-Isberner and Huvermann 2017.  Kind: a survey of
formal results of others.  The chapter prints neither the postulates nor the
relation it asserts between them and ranking independence, and neither is
stated here.  It does print the ranking notion of independence of one
proposition from another, on p.~340: \(B\) is independent of \(A\) when
neither \(A\) nor \(\neg A\) is a reason for \(B\).  Section~\ref{sec:p15-w02b-p15-039-independence}
below proves that this notion is the image of probabilistic independence
under the chapter's translation of products into sums, and that it is
exactly the condition under which no update on \(A\) changes the ranks of
\(B\) and \(\neg B\).  That result is the formal side of the chapter's
remark that can be checked from its own definitions; it is not asserted to
coincide with any postulate of the cited studies
\parencite[pp.~340, 342]{KernIsbernerSkovgaardOlsenSpohn2021Ranking}.

\paragraph{Levi and Harper identities and a common method for revision and contraction (p.~342).}
The chapter says that the main AGM operations, revision and contraction, are
related by the Levi and Harper identities; that in ranking theory the
connection goes deeper, because iterated contraction is a special kind of
iterated ranking conditionalization; and that Kern-Isberner, Bock, Sauerwald
and Beierle 2017 show that iterated revision and contraction can be carried
out by one common method.  Kind: formal.  The identities are those displayed
after Table~\ref{tab:p08-agm-contraction-postulates}.
Theorem~\ref{thm:p15-w02b-p15-039-levi-harper} below proves both for the
chapter's own two-cell update and shows that the contraction it defines is a
case of that update, hence again a ranking state available for further
change.  The 2017 method is described in the chapter only through its
reference title, which names iterated contraction of propositions and
conditionals under the principle of conditional preservation; its
definitions are not printed in the chapter and are not stated here.  The
theorem concerns the two-cell update only
\parencite[pp.~342, 344]{KernIsbernerSkovgaardOlsenSpohn2021Ranking}.

\paragraph{The Bayesian boom (p.~342).}
The chapter reports the observation of Colombo, Elkin and Hartmann 2018 that
probabilistic Bayesian approaches currently attract a great deal of work in
cognitive science, to the neglect of alternative formal frameworks such as
ranking theory.  Kind: a descriptive claim about a research field, reported
at second hand.  The chapter gives no count, period, sampling of the
literature or measure of neglect.  The claim concerns where research
attention goes, not the merit of any framework; nothing about which
framework should be used follows from it, and no such conclusion is drawn
here \parencite[p.~342]{KernIsbernerSkovgaardOlsenSpohn2021Ranking}.

\paragraph{Empirical tests of possibility theory (p.~342).}
The chapter says that Da Silva Neves, Bonnefon and Raufaste 2002 and
Benferhat, Bonnefon and Da Silva Neves 2005 subjected possibility theory to
empirical testing, and it describes only the direct route of the 2002 study
(see Theorem~\ref{thm:p15-w02b-p15-039-rational-monotony} and the paragraph
that follows it).  The 2005 paper is, by its reference title, an overview of
the possibilistic treatment of default reasoning together with experimental
studies; the chapter gives no separate design or result for it.  Kind:
empirical, attributed.  For neither paper does the chapter report task
materials, population, sample size, measures, uncertainty or competing
explanations, and none is supplied here.  Because possibility measures and
ranks share the static ranked structure of
Theorem~\ref{thm:p15-rat53-possibility}, such tests bear on that structure;
the chapter itself stresses that they used possibility judgments and that no
direct test with judgments of disbelief has been made
\parencite[pp.~342, 344]{KernIsbernerSkovgaardOlsenSpohn2021Ranking}.

\paragraph{Language comprehension and background knowledge (pp.~342--343).}
Isberner and Kern-Isberner 2017 are reported to have asked whether revision
with ranking functions retrodicts the delay in processing implausible
information known as the epistemic Stroop effect, under the guiding
assumption that a ranking represents the situation model built during
comprehension.  Eichhorn, Kern-Isberner and Ragni 2018 are reported to have
proposed a conditional-logical model based on ranking functions in which
plausible background knowledge can be elaborated.  Kind: modelling proposals
and a retrodiction, attributed.  The chapter prints neither model's
definitions, inference relation, data or fit; none is stated here.  The
second proposal addresses the gap identified in the suppression discussion of
Chapter~\ref{sec:p15-rationality-5-3}: according to the chapter's account of Ragni and colleagues 2017, the
ranking inference mechanism does not by itself retrodict suppression, and
further assumptions about the knowledge base are required
\parencite[pp.~342--343]{KernIsbernerSkovgaardOlsenSpohn2021Ranking}.

\paragraph{The remaining attributions.}
The antecedents Shackle 1949, Rescher 1964, Cohen 1970, 1977 and 1980 and
Zadeh 1978, and the AI sources Makinson 1989 and Pearl 1990 (p.~341), are
treated in the sections of Chapter~\ref{sec:p15-rationality-5-3}
on comparisons and on AI methods.  The belief
revision research of Baratgin and Politzer 2010, Oaksford and Chater 2013 and
Wolf, Rieger and Knauff 2012, and the principle of conditional preservation
of Kern-Isberner 2004 (p.~342), are treated in the section of that chapter on AI methods and
psychological standing.  Byrne 1989, Ragni and colleagues 2017, Spohn 2013 and
2015, Skovgaard-Olsen 2014, the 2016, 2017 and 2019 reason-relation studies
and Juslin and colleagues 2011 (p.~343) are treated in the same section, the
2016 and 2011 studies from their own pages.  The chapter's further statement
that ranking arithmetic needs much less computational effort than probability
arithmetic is a claim about the elementary operations; it is not a general
complexity comparison, as remarked in that section
\parencite[pp.~341--343]{KernIsbernerSkovgaardOlsenSpohn2021Ranking}.

\section{Revision, contraction and the Levi and Harper identities}
\label{sec:p15-w02b-p15-039-levi-harper}

\begin{hbdefinition}{def:p15-w02b-p15-039-contraction}{Revision, contraction and belief sets for a finite regular ranking}
Let \(W\) be finite and let \(\kappa\) be a negative ranking function whose
world ranks \(r(w)\) are all finite; call such a ranking regular.  For a
regular \(\lambda\) write \(\mathcal B_0(\lambda)\) and \(M_0(\lambda)\) for
the threshold-zero belief set of Definition~\ref{def:p15-rat53-ranking} and
the zero-rank set \(M_0\) of Theorem~\ref{thm:p15-rat53-rank-laws}, both
computed for \(\lambda\).  For nonempty \(A\) let
\(\operatorname{Min}_\kappa(A)=\{w\in A:r(w)=\kappa(A)\}\) be its
minimum-rank worlds.  For a set \(\mathcal K\) of propositions and a
proposition \(E\), the expansion is
\[
 \mathcal K+E=\{A\subseteq W:\ A\supseteq B\land E\ \text{for some}\ B\in\mathcal K\},
\]
the set form of expansion in Definition~\ref{def:p08-agm}.  Let
\(\varnothing\ne E\ne W\) and let \(n\geq1\) be an integer.  Revision by
\(E\) with firmness \(n\) is the update \(\kappa^{0,n}_E\) of
Definition~\ref{def:p15-rat53-soft-update}.  The contraction of \(E\) is
\[
 \kappa^{-}_E=
 \begin{cases}
 \kappa,&\kappa(\neg E)=0,\\
 \kappa^{0,0}_E,&\kappa(\neg E)>0.
 \end{cases}
\]
Regularity and \(\varnothing\ne E\ne W\) make both cell ranks of \(E\)
finite, so both updates are defined.  The second case is the belief-contraction
case of Definition~\ref{def:p15-rat53-soft-update}; the first leaves an
unbelieved \(E\) untouched, as the contraction vacuity postulate of
Table~\ref{tab:p08-agm-contraction-postulates} requires.  The chapter does
not print this case distinction; it is the reading of its statement that
contraction is a special kind of ranking conditionalization adopted here
\parencite[pp.~340, 342]{KernIsbernerSkovgaardOlsenSpohn2021Ranking}.
\end{hbdefinition}

\begin{hbtheorem}{thm:p15-w02b-p15-039-levi-harper}{proved=annex:p15-w02b-p15-039-levi}{The Levi and Harper identities for ranking revision and contraction}
Let \(\kappa\) be regular on finite \(W\), let \(\varnothing\ne E\ne W\)
and let \(n\geq1\).
\begin{enumerate}
\item For every regular \(\lambda\),
 \(\mathcal B_0(\lambda)=\{A:M_0(\lambda)\subseteq A\}\) and
 \(\mathcal B_0(\lambda)+E=\{A:M_0(\lambda)\land E\subseteq A\}\).
 Moreover \(M_0(\kappa^{0,n}_E)=\operatorname{Min}_\kappa(E)\) and
 \(M_0(\kappa^{0,0}_E)=\operatorname{Min}_\kappa(E)\lor\operatorname{Min}_\kappa(\neg E)\).
\item \(\kappa^{-}_E\) is regular,
 \(\mathcal B_0(\kappa^{-}_E)\subseteq\mathcal B_0(\kappa)\) and
 \(E\notin\mathcal B_0(\kappa^{-}_E)\).
\item Levi identity:
 \(\mathcal B_0(\kappa^{0,n}_E)=\mathcal B_0(\kappa^{-}_{\neg E})+E\).
\item Harper identity:
 \(\mathcal B_0(\kappa^{-}_E)=\mathcal B_0(\kappa)\cap\mathcal B_0(\kappa^{0,n}_{\neg E})\).
\item On ranking states,
 \((\kappa^{-}_{\neg E})^{0,n}_E=\kappa^{0,n}_E\).
\end{enumerate}
\end{hbtheorem}

Parts~(iii) and~(iv) are the two identities of
Definition~\ref{def:p08-agm} with \(\mathcal B_0\) in place of a belief set.
Part~(v) is stronger than~(iii): contracting the negation and then revising
by \(E\) returns the very ranking state that direct revision gives, not only
its belief set.  Part~(ii) gives inclusion and success; it also shows that a
contraction leaves a regular ranking, so that a further contraction or
revision is defined.  This is the precise sense in which iterated contraction
is iterated conditionalization here.  The theorem does not cover \(E=W\),
\(E=\varnothing\) or a ranking with an infinite world rank: in those cases a
cell of \(E\) can have infinite rank, and the subtraction rule of
Definition~\ref{def:p15-rat53-conditionals} does not apply.

\begin{hbexample}{def:p15-w02b-p15-039-levi-example}{Contraction in the two rankings of the revision experiments}
Take the rankings \(\kappa_1,\kappa_2\) displayed in
Chapter~\ref{sec:p15-rationality-5-3} for the belief revision experiments,
with \(E=\neg q\) and \(n=1\).  Both have \(\kappa_i(q)=0\) and
\(\kappa_i(\neg q)=1\), so \(q\) is believed and its contraction is
\((\kappa_i)^{0,0}_q\): the \(q\)-worlds keep their ranks and each
\(\neg q\)-world drops by one.

For \(\kappa_1\) the result ranks \((1,1)\) and \((1,0)\) at \(0\),
\((0,0)\) at \(1\) and \((0,1)\) at \(3\).  Its zero set consists of the two
\(p\)-worlds, so the contracted belief set is that of \(p\).  Expanding by
\(\neg q\) leaves \((1,0)\), which is \(\operatorname{Min}_{\kappa_1}(\neg q)\),
the zero set of the revision \((\kappa_1)^{0,1}_{\neg q}\): the Levi identity
in this instance.  Conversely the zero sets of \(\kappa_1\) and of that
revision are \(\{(1,1)\}\) and \(\{(1,0)\}\); their union is the zero set of
the contraction: the Harper identity.

For \(\kappa_2\) the result ranks \((1,1)\) and \((0,0)\) at \(0\),
\((1,0)\) at \(1\) and \((0,1)\) at \(3\), and the contracted belief set is
that of \(p\leftrightarrow q\), which keeps \(p\to q\) and gives up \(p\).
Thus two states with the same belief set
\(\operatorname{Cn}(\{p\to q,p\})\) contract \(q\) differently.  The
identities pass between revision and contraction for one state, but neither
operation is a function of the belief set alone; the ranks decide, exactly as
for the revisions of Theorem~\ref{thm:p19-r54-two-revisions}.
\end{hbexample}

\section{Independence of propositions and revision}
\label{sec:p15-w02b-p15-039-independence}

The chapter defines independence through its reason relation: \(B\) is
irrelevant to, or independent of, \(A\) when neither \(A\) nor \(\neg A\) is
a reason for \(B\).  It adds that a theory of independence and conditional
independence can be built in ranking terms in wide analogy with the
probabilistic theory, and on p.~342 it relates ranking independence to the
independence properties proposed for iterated and multiple revision.  The
next result fixes what the definition amounts to for one pair of
propositions, using the cell ranks
\(a=\kappa(A\land B)\), \(b=\kappa(A\land\neg B)\), \(c=\kappa(\neg A\land B)\),
\(d=\kappa(\neg A\land\neg B)\) of
Theorem~\ref{thm:p15-w02b-p15-039-reason-laws} and the update
\(\kappa^{x,y}_A\) of Definition~\ref{def:p15-rat53-soft-update}, whose new
cell ranks are written \(x,y\) here to keep them apart from the cell ranks
\parencite[pp.~340, 342]{KernIsbernerSkovgaardOlsenSpohn2021Ranking}.

\begin{hbtheorem}{thm:p15-w02b-p15-039-independence}{proved=annex:p15-w02b-p15-039-independence}{Ranking independence: product form and invariance under update}
Let \(\kappa\) be a negative ranking function and let the four cell ranks
\(a,b,c,d\) be finite.  The following are equivalent.
\begin{enumerate}
\item \(B\) is independent of \(A\): neither \(A\) nor \(\neg A\) is a
 reason for \(B\).
\item \(\tau(B\mid A)=\tau(B\mid\neg A)\).
\item \(a+d=b+c\).
\item \(\kappa(A'\land B')=\kappa(A')+\kappa(B')\) for every
 \(A'\in\{A,\neg A\}\) and \(B'\in\{B,\neg B\}\).
\item For all \(x,y\in\mathbb N\) with \(\min\{x,y\}=0\),
 \(\kappa^{x,y}_A(B)=\kappa(B)\) and \(\kappa^{x,y}_A(\neg B)=\kappa(\neg B)\).
\end{enumerate}
Moreover, independence is symmetric: \(B\) is independent of \(A\) exactly
when \(A\) is independent of \(B\).  And every update \(\kappa^{x,y}_A\)
leaves \(\tau(B\mid A)\) and \(\tau(B\mid\neg A)\) unchanged, so it
preserves whether \(A\) is a reason for \(B\), a reason against \(B\), or
neither.
\end{hbtheorem}

Item~4 is the translation, announced by the chapter on p.~341, of
\(P(A\land B)=P(A)P(B)\) in Definition~\ref{def:p15-finite-independence}:
a product of probabilities becomes a sum of negative ranks.  Item~4 states
the sum equation for all four cell pairs, and the proof shows that the single
cross-sum equation of item~3 yields all four at once.  Item~5 is a property of
revision.  Whatever is learned about \(A\), with whatever firmness, the
ranks of \(B\) and of \(\neg B\) stay fixed exactly when \(B\) is
independent of \(A\); so belief, disbelief or suspension about \(B\) and
its firmness survive every update on \(A\).  This is the kind of property
that the chapter connects with the independence postulates of the revision
studies it cites.  Whether it coincides with any of those postulates cannot
be read off the chapter, which does not print them.

The theorem concerns two propositions.  It does not give conditional
independence given a third proposition, the independence of variables, or
the ranking version of the Bayes-net theory that the chapter attributes to
Spohn 1988 and to Goldszmidt and Pearl 1996; those are not derived here.
Nor is independence preserved by every change.  An update on \(A\) keeps it,
by the last statement of the theorem, but an update on another proposition
need not, as the following example shows.

\begin{hbexample}{def:p15-w02b-p15-039-independence-example}{Independence, its preservation, and its loss}
Let \(W\) consist of four worlds, one in each cell, with
\((a,b,c,d)=(0,1,2,3)\).  Then \(a+d=3=b+c\), so \(B\) is independent of
\(A\).  The marginals are \(\kappa(A)=0\), \(\kappa(\neg A)=2\),
\(\kappa(B)=0\), \(\kappa(\neg B)=1\), and each cell is the sum of its two
marginals, for instance \(d=3=2+1\).  Accepting \(\neg A\) with firmness
\(1\), the update \(\kappa^{1,0}_A\), gives cells \((1,2,0,1)\): now
\(\neg A\) is believed, the cross sums are still \(2\) and \(2\), and
\(\kappa(B)=0\), \(\kappa(\neg B)=1\) as before.  Accepting instead the
proposition \(E=(A\land B)\lor(\neg A\land\neg B)\), which is already
believed with \(\kappa(E)=0\) and \(\kappa(\neg E)=1\), with firmness \(2\),
the update \(\kappa^{0,2}_E\), keeps the ranks of the \(E\)-worlds and moves
each \(\neg E\)-world to \(2+(r(w)-\kappa(\neg E))=r(w)+1\), giving cells
\((0,2,3,3)\).  The cross sums are now \(3\) and \(5\): \(A\) has become a
reason for \(B\).

The Tweety prior of Table~\ref{tab:p15-rat53-tweety} shows the last
statement of the theorem with a dependent pair.  Its cells for bird and
penguin are \(\kappa(B\land P)=1\), \(\kappa(B\land\neg P)=0\),
\(\kappa(\neg B\land P)=8\), \(\kappa(\neg B\land\neg P)=0\), with cross sums
\(1\) and \(8\), so being a bird is a reason for being a penguin.  The bird
evidence of the table is an update on \(B\); afterwards the cells are
\(1,0,10,2\), and \(\tau(P\mid B)=0-1=-1\) and
\(\tau(P\mid\neg B)=2-10=-8\) are the prior values \(0-1\) and \(0-8\).
\end{hbexample}

\section{Consonant belief functions and ranks}
\label{sec:p15-w02b-p15-039-consonant}

The functions \(m\), \(\operatorname{Bel}\), \(\operatorname{Pl}\) and the
normalized combination \(\otimes\) are those of
Definitions~\ref{def:p15-r47-mass} and~\ref{def:p15-r47-combination}, with
the finite space \(W\) in place of \(S\); \(\Pi\) and \(N\) are the
transforms of Theorem~\ref{thm:p15-rat53-possibility}.

\begin{hbtheorem}{thm:p15-w02b-p15-039-consonant}{proved=annex:p15-w02b-p15-039-consonant}{Ranks as consonant belief functions: agreement on conditioning, divergence under combination}
Let \(W\) be finite and let \(\kappa\) be a negative ranking function with
world ranks \(r\).  Let \(0=k_0<k_1<\cdots<k_t\) be the distinct finite
values of \(r\) and \(F_j=\{w:r(w)\leq k_j\}\).  Put
\(m_\kappa(F_j)=2^{-k_j}-2^{-k_{j+1}}\) for \(j<t\),
\(m_\kappa(F_t)=2^{-k_t}\), and \(m_\kappa=0\) on every other set.
\begin{enumerate}
\item \(m_\kappa\) is a basic assignment whose focal sets
 \(F_0\subset F_1\subset\cdots\subset F_t\) form a chain, and for every
 \(A\subseteq W\)
 \[
  \operatorname{Pl}(A)=2^{-\kappa(A)}=\Pi(A),\qquad
  \operatorname{Bel}(A)=1-2^{-\kappa(\neg A)}=N(A).
 \]
\item If \(\kappa(E)<\infty\) and \(m^E\) places mass one on \(E\), then
 \(m_\kappa\) and \(m^E\) have \(K=2^{-\kappa(E)}>0\), the focal sets of
 \(m_\kappa\otimes m^E\) form a chain, and its plausibility is
 \(A\mapsto2^{-\kappa(A\mid E)}\), with \(2^{-\infty}=0\).
\item On \(W=\{x,y,z\}\) let \(\lambda_1\) and \(\lambda_2\) have world
 ranks \((0,0,1)\) and \((1,0,0)\) at \((x,y,z)\).  Then
 \(m_{\lambda_1}\otimes m_{\lambda_2}\) gives mass \(1/4\) to each of
 \(\{y\}\), \(\{x,y\}\), \(\{y,z\}\) and \(W\); its belief has
 \(\operatorname{Bel}(\{x,y\})=\operatorname{Bel}(\{y,z\})=1/2\) and
 \(\operatorname{Bel}(\{y\})=1/4\), and it equals \(m_\mu\) for no negative
 ranking function \(\mu\) on \(W\).
\end{enumerate}
\end{hbtheorem}

Part~(i) makes the chapter's static equivalence exact on a finite space:
every ranking is represented by a consonant belief function through the
transform \(2^{-\kappa}\).  Part~(ii) shows that conditioning agrees: adding
the certain information \(E\) by Dempster's rule produces the consonant
function of the conditional ranking \(\kappa(\cdot\mid E)\).  Part~(iii) shows
that the agreement ends with combination: two pieces of information, each
represented by a ranking, are combined by Dempster's rule into a belief
function whose focal sets do not form a chain and whose belief is not
minitive, so it represents no ranking.  The chapter's ranking dynamics, the
update of Definition~\ref{def:p15-rat53-soft-update}, is specified by new
cell ranks and is not a combination rule; part~(iii) therefore does not say
what ranking theory puts in place of Dempster's rule, only that the static
correspondence is not preserved by that rule.  The transform is fixed with
base \(2\); another base gives a different consonant function with the same
order, so the numerical masses are a choice of representation, not a
measurement of disbelief.
